% Judul Capstone
\judul{Pembuatan \textit{Template} Laporan \textit{Capstone Project} dalam Format LaTeX}
\title{Template Design for Capstone Project Report in LaTeX Format}

% Jenis Dokumen
%% Jenis Dokumen : PERANCANGAN PRODUK DAN SPESIFIKASI

% Kode Dokumen
%% Kode Dokumen : C-251

% Nomor Dokumen (ID Kelompok Capstone)
\NoDok{C\_06}

% Nomor Revisi
\NoRev{02}

% Tanggal Penerbitan Dokumen
%% Otomatis terisi tanggal ketika file LaTeX ini di-compile

% Data Mahasiswa Capstone
%% Format : \MHS{<Nama Lengkap>}{<NIM>}{<Prodi>}{<Alamat Email>}
% Ketua Kelompok
\MHSA{Irsad Najib Eka Putra}{23/518119/TK/57005}
	  {Tekonologi Informasi}{irsadnajibekaputra@mail.ugm.ac.id}
% Anggota 1
\MHSB{Nafil Nissano Yogma Pratama}{23/521136/TK/57475}
	  {Teknik Biomedis}{nafilnissanoyogmapra@mail.ugm.ac.id} 
% Anggota 2
\MHSC{Azizah Az Zahra}{23/519983/TK/57292}
	  {Teknologi Biomedis}{azizahazzahra2003@mail.ugm.ac.id}
% Anggota 3
\MHSD{Isna Rahmanadia}{23/520928/TK/57450}
	  {Teknologi Biomedis}{isnarahmanadia2005@mail.ugm.ac.id}
% Anggota 4
\MHSE{Dimitris G. Manolakis}{20/443100/TK/54200}
	  {Teknik Elektro}{dgmanolakis@mail.ugm.ac.id}
% Un-comment Line di bawah ini apabila tidak ada Anggota 4 :
\MHSE{}{}{}{}

% Dosen Pembimbing
%% Format : {<Nama Lengkap>}{<NIP/NIU>}
\DPA{Carl Friedrich Gauss, B.Eng., M.Eng., D.Eng.}{111 1993 01 2024 01 101}

% Tempat Pelaksanaan
%% Format : {<Nama Laboratorium> \newline <Nama Departemen> \newline <Nama Fakultas>}
\Tempat{Departemen Teknik Elektro dan Teknologi Informasi \newline Fakultas Teknik}